\section{Optimización y Despliegue en el Borde}

Los modelos empleados en el sistema son de gran porte ---la variante de mayor
tamaño supera los cien megabytes en sus pesos originales---, mientras que el
destino de ejecución es un equipo con recursos acotados que debe procesar video
de forma continua. Este apartado documenta las transformaciones aplicadas a los
modelos para reducir su costo de ejecución y las estrategias que disminuyen el
volumen de cómputo requerido, y cierra señalando qué mediciones restan para
cuantificar su efecto.

\subsection{Exportación a formatos intermedios}

Los pesos originales de la familia de detectores se distribuyen en un formato
orientado al entrenamiento, que no es el más eficiente para inferencia. El
primer paso consiste en exportarlos a un formato intermedio de grafo de cómputo
que describa la red de manera independiente del marco de trabajo, lo que permite
desplegarla sobre distintos ejecutores.

La exportación se realiza con una resolución de entrada fija, condición necesaria
para que el compilador posterior pueda especializar las operaciones. Cuando el
formato admite el almacenamiento de los pesos como datos externos, el grafo y los
parámetros quedan en archivos separados; la Tabla \ref{tab:formatos} registra la
huella efectiva de cada formato para la variante empleada en la extracción.

\begin{longtable}{
        P{0.18}
        P{0.12}
        P{0.15}
        P{0.45}
    }
    \caption{Huella en disco por formato de modelo}
    \label{tab:formatos}\\

    \toprule
    \rowcolor{gray!20}
    \textbf{Formato} &
    \textbf{Precisión} &
    \textbf{Tamaño} &
    \textbf{Observación} \\
    \midrule
    \endfirsthead

    \multicolumn{4}{c}{\textit{Continuación de la Tabla
    \ref{tab:formatos}}} \\
    \toprule
    \rowcolor{gray!20}
    \textbf{Formato} &
    \textbf{Precisión} &
    \textbf{Tamaño} &
    \textbf{Observación} \\
    \midrule
    \endhead

    Pesos (variante pequeña) & FP32 & $4{.}8$ MB & Alternativa más liviana. \\
    \midrule
    Pesos (variante media) & FP32 & $42{.}2$ MB & Alternativa intermedia. \\
    \midrule
    Pesos (variante grande) & FP32 & $113{.}2$ MB & Variante empleada en la
    extracción. \\
    \midrule
    Grafo (FP32) & FP32 & $212{.}9$ MB $+$ $214{.}4$ MB &
    Pesos almacenados como datos externos. \\
    \midrule
    Grafo (FP16) & FP16 & $106{.}6$ MB & Media precisión. \\
    \midrule
    Motor compilado & FP16 & $345{.}5$ MB & Específico del dispositivo de
    compilación. \\*
    \multicolumn{4}{r}{\footnotesize\textit{Fuente: Elaboración propia}} \\*
    \bottomrule
\end{longtable}

La huella del artefacto en media precisión ($106{.}6$ MB) frente a la del grafo en
precisión simple con pesos externos ($212{.}9$ MB $+$ $214{.}4$ MB $=$ $427{.}3$ MB)
representa una reducción de aproximadamente el $75\%$. Esa diferencia no proviene
únicamente del cambio de precisión: FP16, por sí solo, reduciría los pesos a la
mitad. La cifra corresponde a la huella efectiva del artefacto exportado y
simplificado, de modo que se reporta como reducción del artefacto y no como efecto
aislado de la media precisión. El motor compilado, en cambio, ocupa un espacio
mayor que los pesos originales: un motor de inferencia contiene, además de los
parámetros, los núcleos de ejecución y las estrategias de paralelización
seleccionadas para el dispositivo. Se trata, por tanto, de un intercambio
deliberado de espacio en disco por velocidad de ejecución, cuyo efecto sobre la
latencia no se midió en este trabajo.

\subsection{Compilación y ejecución con TensorRT}

El formato intermedio se compila finalmente a un motor especializado para el
acelerador gráfico de destino. La compilación toma el grafo exportado y la
resolución de entrada, y genera un artefacto optimizado que incluye las
operaciones fusionadas y los núcleos seleccionados para la arquitectura
específica del dispositivo.

La selección del motor se realiza de forma automática: la configuración resuelve,
para cada modelo, la existencia de un motor compilado junto a los pesos y lo
prefiere cuando está disponible, recurriendo a los pesos originales en los
entornos donde el motor no puede ejecutarse. Gracias a esta regla, la cadena de
herramientas no requiere modificación alguna al pasar de un entorno de cómputo a
otro: basta con compilar el motor una vez en el equipo de destino.

Cabe precisar una limitación relevante para el despliegue: \textbf{los motores
compilados son específicos del dispositivo y del controlador} con los que fueron
generados. Un motor construido para un acelerador determinado no es portable a
otro, por lo que cada equipo de destino requiere su propia compilación. Esta
característica condiciona la distribución del sistema y se declara en el apartado
de limitaciones.

En cuanto a la precisión de ejecución, el motor se compila en media precisión y
los pesos que se ejecutan sin compilación se cargan en ese mismo modo. Los
motores ya incorporan la precisión con la que fueron compilados, de modo que no
se aplica una conversión adicional al cargarlos.

\subsection{Estrategias de reducción de carga computacional}

Más allá de la transformación de los modelos, la arquitectura del proceso reduce
el cómputo requerido mediante decisiones de diseño que evitan trabajo
innecesario. Las estrategias aplicadas son las siguientes:

\begin{enumerate}
    \item \textbf{Muestreo disperso previo.} La primera pasada del extractor
        examina una de cada doce muestras del video con una señal de movimiento
        de costo casi nulo, de modo que la ejecución de la red se reserva para
        los intervalos con actividad.

    \item \textbf{Procesamiento restringido a ventanas de actividad.} La segunda
        pasada decodifica únicamente los intervalos marcados, en lugar del video
        completo. La reducción es proporcional a la fracción de tiempo con
        vehículos presentes.

    \item \textbf{Enmascarado de la información incrustada.} Las regiones donde el
        grabador superpone la etiqueta de canal y la marca temporal se anulan
        antes de invocar al reconocimiento de texto, evitando un costo de
        inferencia que además produciría lecturas espurias.

    \item \textbf{Selección de los mejores fotogramas por trayectoria.} Sólo los
        fotogramas mejor puntuados de cada vehículo alcanzan la etapa de
        reconocimiento, en lugar de la totalidad de los fotogramas en los que
        aparece.

    \item \textbf{Filtros de calidad previos al reconocimiento.} Se exige un
        tamaño mínimo y una nitidez mínima al recorte de placa, y se descartan
        las detecciones de área inferior al umbral establecido. Estos filtros
        evitan invocar al motor de reconocimiento sobre recortes sin posibilidad
        de lectura.

    \item \textbf{Ejecución asíncrona con concurrencia unitaria.} El trabajador de
        inferencia se ejecuta con una sola unidad de concurrencia, dado que cada
        tarea satura el acelerador. Aumentar la concurrencia no reduciría la
        latencia y, en cambio, competiría por la misma unidad de cómputo.

    \item \textbf{Media precisión.} Tanto el motor compilado como los pesos que se
        ejecutan sin compilación operan en media precisión.
\end{enumerate}

El sistema registra, desde la implementación del trabajador, la latencia de
inferencia de cada fotograma procesado: cada observación persiste el tiempo que
demandó el motor de visión para ese fotograma, de modo que la base de datos ya
acumula el dato que el despliegue necesita. Lo que resta es el banco comparativo
que ese registro no reemplaza: la misma secuencia procesada con cada formato de
modelo y cada variante de escala, sobre el hardware de destino, junto con la
degradación de exactitud atribuible a la media precisión. Esa comparación es la
que permite decidir si el intercambio entre velocidad y precisión es aceptable
para un punto de acceso concreto, y no se realizó en este trabajo.
